\documentclass[11pt]{article}
\usepackage[a4paper, margin=1in]{geometry}
\usepackage[T1]{fontenc}
\usepackage[utf8]{inputenc}
\usepackage{amsmath, amssymb}
\usepackage{booktabs}
\usepackage{hyperref}
\usepackage{xcolor}
\hypersetup{colorlinks=true, linkcolor=blue, urlcolor=blue, citecolor=blue}
\usepackage{authblk}
\usepackage{microtype}
\usepackage{enumitem}
\setlist[itemize]{leftmargin=1.4em}

\title{\textbf{RealQM and the GAD65 Autoantigen in Type-1 Diabetes}\\[0.3em]
\large From first-principles PLP electronics to the apo--holo transition and autoimmunity}

\author[1]{Claes Johnson}
\author[2]{\r{A}ke Lernmark}
\affil[1]{KTH Royal Institute of Technology, Stockholm, Sweden}
\affil[2]{Lund University Diabetes Centre, Department of Clinical Sciences, Lund University, Malm\"o, Sweden}
\date{\today \\[0.5em]
\small\fbox{\parbox{0.86\textwidth}{\centering \textbf{DRAFT for co-author review.} The type-1-diabetes
clinical and autoimmunity content is written from the established literature and is to be finalised and
verified by the co-author; the computational (RealQM) results are the present author's. Prepared in dialogue
with Claude (Anthropic).}}}

\begin{document}
\maketitle

\begin{abstract}
Glutamic acid decarboxylase GAD65 is both an enzyme --- it makes the inhibitory neurotransmitter GABA from
L-glutamate using the cofactor pyridoxal-5$'$-phosphate (PLP) --- and the \textbf{major autoantigen of
type-1 diabetes (T1D)}: autoantibodies to GAD65 (GADA) are a front-line predictive and diagnostic marker.
GAD65's autoantigenicity is linked to its unusual catalytic behaviour: unlike its stable paralogue GAD67 it
readily \emph{loses PLP} and populates a flexible \textbf{apo} form, and its mobile catalytic loop overlaps
conformational GADA epitopes. Any mechanistic account of this chemistry--immunology link needs the underlying
PLP electronics right. We bring to bear \textbf{RealQM}, a parameter-free, real-space electronic-structure
method (no basis set, no exchange--correlation functional) cheap enough to \emph{scan}. We show that RealQM,
given only the active-site nuclei, \emph{spontaneously} reproduces the PLP \textbf{electron-sink} mechanism
that defines GAD65 catalysis --- with an identical-substrate control isolating the conjugated ring as the
cause --- and we propose a joint program that carries this first-principles handle from the catalytic step to
the \textbf{holo$\leftrightarrow$apo} transition and the catalytic-loop/epitope region, i.e.\ toward a
computational account of \emph{why GAD65 is the autoantigen}. This is a proposal and a first result, not a
finished clinical claim.
\end{abstract}

\section{GAD65: enzyme and type-1-diabetes autoantigen}
Type-1 diabetes is a T-cell-mediated autoimmune destruction of the insulin-producing pancreatic $\beta$-cells,
preceded for months to years by circulating islet autoantibodies. Among these, autoantibodies to the 65-kDa
isoform of glutamic acid decarboxylase --- \textbf{GADA} --- are one of the earliest and most informative:
GAD65 was identified as the ``64K'' islet autoantigen of insulin-dependent diabetes~\cite{Baekkeskov1990}, and
GADA is now a standard marker for risk prediction, staging, and classification (including LADA in
adults)~\cite{Lernmark_review}. Enzymatically, GAD65 is a PLP-dependent decarboxylase that converts
L-glutamate to GABA + CO$_2$.

Two features make GAD65, rather than the closely related and more catalytically stable GAD67, the dominant
autoantigen. First, GAD65 undergoes \textbf{autoinactivation}: it releases its PLP cofactor and cycles between
a catalytically competent \emph{holo} form and a flexible, apo form; the structural basis is a mobile
``catalytic loop'' region resolved crystallographically~\cite{Fenalti2007}. Second, the principal GADA
epitopes are \textbf{conformational} and cluster in the middle/C-terminal regions whose flexibility is tied to
this apo$\leftrightarrow$holo cycling. The prevailing view therefore connects GAD65's \emph{chemistry} (PLP
handling, conformational dynamics) to its \emph{immunogenicity} (epitope exposure). Testing that connection
mechanistically is the target of this collaboration.

\section{The autoantigenicity puzzle, and where computation can help}
The chemistry--immunology link above is qualitative. A quantitative, first-principles account would need to
answer: how does the active-site electronic structure differ between holo and apo GAD65? Does PLP loss (or an
altered PLP electronic state) drive the loop conformational changes that expose conformational epitopes? Why
GAD65 and not GAD67? These are electronic-structure questions at heart, but at a scale (the full flexible
enzyme, many conformations) that the standard first-principles tools reach only with difficulty. We propose
that a \emph{cheap, parameter-free} electronic-structure method, validated first on the catalytic step, is the
right instrument to open this regime.

\section{RealQM: a parameter-free electronic handle}
RealQM formulates quantum mechanics as real-space, grid-based continuum mechanics: unit-charge electron
densities relax to minimum Coulomb energy in the field of fixed nuclei, with \textbf{no basis set and no
exchange--correlation functional}~\cite{realqm}. Two properties matter here. (i) It is \emph{parameter-free}
--- there is no functional whose choice can change the answer, DFT's principal systematic-error source. (ii)
Its cost is set by the grid rather than by $N^3$ in the electron number, so it is cheap enough to \emph{scan}
many configurations on commodity hardware --- the throughput a holo-vs-apo and loop-conformation study
requires. Its accuracy is qualitative on absolute energies but good on relative, structural observables (it
reproduces, e.g., the water dipole to ${\sim}99\%$); the appropriate use here is trend and mechanism, not
calibrated barriers, for which mature QM/MM remains the reference.

\section{Result: the PLP electron sink from first principles}
In PLP decarboxylases the scissile C$\alpha$--COO$^-$ bond breaks with the departing electrons delocalising
into the protonated pyridine ring --- the quinonoid intermediate --- with the pyridinium N the
\textbf{electron sink}. We tested whether RealQM reproduces this given only nuclei, on a reduced gas-phase
active-site model (protonated pyridine sink + 3-OH + 4$'$-imine; glutamate external aldimine, Dunathan-aligned).
The breaking C$\alpha$--CO$_2$ distance $R$ was driven outward, electrons relaxed at each step, for two
environments: glutamate \textbf{conjugated} to the ring (sink present) and the \emph{identical} glutamate as a
free amino acid (no sink --- a built-in control). The sink term is $\Delta E_{\text{free}}-\Delta E_{\text{PLP}}$.
\begin{center}
\begin{tabular}{rrrr}
\toprule
$R$ (\AA) & $\Delta E$ free & $\Delta E$ PLP & sink \\
\midrule
1.55 & $0.000$ & $0.000$ & $0.00$ \\
2.60 & $-2.476$ & $-3.776$ & $+1.30$ \\
4.20 & $-3.101$ & $-6.605$ & $+3.50$ \\
6.00 & $-3.454$ & $\mathbf{-7.620}$ & $\mathbf{+4.17}$ \\
\bottomrule
\end{tabular}
\end{center}
The two curves diverge: free glutamate \emph{saturates} (${\approx}-3.5$ once CO$_2$ is ${\sim}4$~\AA{} out ---
no sink, a dead end), while the conjugated case \emph{keeps descending} (no plateau); the sink term grows
monotonically to $+4.2$ (model units) and is still climbing at $6$~\AA. With no reaction coordinate built in,
RealQM recovers the textbook role of PLP: the conjugated ring, and only it, stabilises the carbanion. This is
a \emph{qualitative} validation (sign/trend, uncalibrated; thermodynamic, not a barrier), but it establishes
that the method captures the central electronic event of GAD65 catalysis from first principles --- the
prerequisite for the harder questions.

\section{The joint program: from catalytic electronics to autoantigenicity}
We propose the following staged program, combining RealQM computation (Johnson) with GAD65/T1D expertise
(Lernmark):
\begin{itemize}
\item \textbf{Aim 1 --- calibrate.} Benchmark RealQM against high-level QM/MM on the GAD65 catalytic step
(quinonoid formation), establishing the trends it can be trusted on.
\item \textbf{Aim 2 --- holo vs apo electronics.} Compare the active-site electronic stabilisation of the
PLP-bound (holo) and cofactor-free (apo) states; test whether the electronic change accompanies/drives the
catalytic-loop rearrangement.
\item \textbf{Aim 3 --- loop/epitope coupling.} Map the electronic and conformational state of the catalytic
region that overlaps conformational GADA epitopes; ask whether apo-state features correlate with epitope
exposure --- and why GAD65 differs from the stable GAD67.
\item \textbf{Aim 4 --- recognition (longer term).} Model GADA--epitope recognition with real, conformational
epitope and CDR residues (a substantial lift beyond the current placeholder docking model).
\end{itemize}
The value is throughput and a functional-free cross-check on a regime --- GAD65 \emph{specifically}, and its
apo/autoantigenic state --- that generic PLP-decarboxylase QM/MM has largely left open.

\section{Clinical relevance and outlook}
If the apo-state electronics/conformation can be connected computationally to conformational-epitope exposure,
the payoff is mechanistic: a first-principles rationale for why GAD65 (not GAD67) is the T1D autoantigen, and
a route toward understanding and perhaps predicting GADA epitope specificity --- of direct interest for risk
stratification and for antigen-specific interventions. The present paper contributes the computational
foundation (a parameter-free method that gets GAD65's central chemistry right) and the joint plan; the clinical
and immunological interpretation is to be developed with the co-author.

\section*{Author contributions}
C.J.\ developed the RealQM methodology and performed the computations. \r{A}.L.\ [to be completed] contributes
the GAD65/type-1-diabetes autoimmunity expertise and clinical interpretation. Both authors will finalise the
manuscript.

\begin{thebibliography}{9}
\bibitem{Baekkeskov1990} S.\ Baekkeskov, H.-J.\ Aanstoot, S.\ Christgau, \emph{et al.} (incl.\ \r{A}.\
Lernmark), \emph{Identification of the 64K autoantigen in insulin-dependent diabetes as the GABA-synthesizing
enzyme glutamic acid decarboxylase}, Nature \textbf{347}, 151--156 (1990).
\bibitem{Fenalti2007} G.\ Fenalti \emph{et al.}, \emph{GABA production by glutamic acid decarboxylase is
regulated by a dynamic catalytic loop}, Nat.\ Struct.\ Mol.\ Biol.\ \textbf{14}, 280--286 (2007).
\bibitem{Lernmark_review} \r{A}.\ Lernmark \emph{et al.}, reviews on GAD65 autoantibodies (GADA) as
predictive and diagnostic markers in type-1 diabetes and LADA. [to be completed with the co-author's preferred
references.]
\bibitem{realqm} C.\ Johnson, \emph{RealQM: Quantum Mechanics as Multiphase 3D Continuum Mechanics} (arXiv
article and interactive gallery), \url{https://claes542.github.io/RealMolecule/}. Companion single-author
paper on the GAD65 electron-sink test therein.
\end{thebibliography}

\end{document}
